% !TeX root = tcolorbox.tex
% include file of tcolorbox.tex (manual of the LaTeX package tcolorbox)
\clearpage
\section{\mylib{hooks} 程序库}\label{sec:hooks}%
\tcbset{external/prefix=external/hooks_}%
该程序库通过宏包选项或在导言区中用以下方式加载：
\begin{dispListing}
\tcbuselibrary{hooks}
\end{dispListing}

对于与皮肤相关的选项，还需单独加载 \mylib{skins} 程序库。


\subsection{钩子的概念}
钩子（hook）是 \LaTeX\ 代码中的一个占位位置，可在该处追加额外代码。
例如，\LaTeX\ 宏 |\AtBeginDocument| 会把代码挂到一个钩子上，
该钩子位于每份文档的开头。

|tcolorbox| 的若干选项键允许提供一段代码，
把它添加到彩色文本框的特定位置。例如，\refKey{/tcb/before upper}
会把代码放在上部内容之前。之后再次使用该键会覆盖此前的设置。

\mylib{hooks} 程序库把 \refKey{/tcb/before upper} 等已有键扩展为
\enquote{可挂接}版本，例如 \refKey{/tcb/before upper app} 和
\refKey{/tcb/before upper pre}。这些\enquote{可挂接}键不会覆盖已有设置，
而是把新给出的代码\emph{追加}（app）或\emph{前插}（pre）到已有代码之后或之前。

一般命名约定（有少数小例外）为：
\begin{itemize}
\item \meta{option key} |app|：作用与 \meta{option key} 相同，但
  把其代码\emph{追加}到已有代码之后。
\item \meta{option key} |pre|：作用与 \meta{option key} 相同，但
  把其代码\emph{前插}到已有代码之前。
\end{itemize}
若（再次）使用原始的 \meta{option key}，则所有代码都会被覆盖。
因此，使用各选项键的先后顺序至关重要。

特别地，若原始 \meta{option key} 本身默认即可堆叠，
通常已有一个配套的\enquote{|no| \meta{option key}}选项用于清空已有代码。
因此本程序库只补充\enquote{\meta{option key} |pre|}这一变体。

\enlargethispage*{2cm}

\begin{dispExample}
% \usepackage{array,tabularx}
\newcolumntype{Y}{>{\raggedleft\arraybackslash}X}% see tabularx
\tcbset{enhanced,fonttitle=\bfseries\large,fontupper=\normalsize\sffamily,
  colback=yellow!10!white,colframe=red!50!black,colbacktitle=Salmon!30!white,
  coltitle=black,center title,
  tabularx={X||Y|Y|Y|Y||Y},% this sets `before upper' and `after upper'
  before upper app={Group & One & Two & Three & Four & Sum\\\hline\hline}  }

\begin{tcolorbox}[title=My table]
Red   & 1000.00 & 2000.00 &  3000.00 &  4000.00 & 10000.00\\\hline
Green & 2000.00 & 3000.00 &  4000.00 &  5000.00 & 14000.00\\\hline
Blue  & 3000.00 & 4000.00 &  5000.00 &  6000.00 & 18000.00\\\hline\hline
Sum   & 6000.00 & 9000.00 & 12000.00 & 15000.00 & 42000.00
\end{tcolorbox}
\end{dispExample}



\clearpage
\subsection{文本框内容的追加}\label{subsec:hookscontentadditions}
以下各选项键是对 \zvref{subsec:contentadditions} 中已有选项的扩展。

\begin{docTcbKey}{before title app}{=\meta{code}}{无默认值}
  把给定的 \meta{code} \emph{追加}到 \refKey{/tcb/before title}，位置在颜色与字体设置\emph{之后}、标题内容\emph{之前}。
\end{docTcbKey}

\begin{docTcbKey}{before title pre}{=\meta{code}}{无默认值}
  把给定的 \meta{code} \emph{前插}到 \refKey{/tcb/before title}，位置在颜色与字体设置\emph{之后}、标题内容\emph{之前}。
\end{docTcbKey}

\begin{docTcbKey}{after title app}{=\meta{code}}{无默认值}
  把给定的 \meta{code} \emph{追加}到 \refKey{/tcb/after title}，位置在标题内容\emph{之后}。
\end{docTcbKey}

\begin{docTcbKey}{after title pre}{=\meta{code}}{无默认值}
  把给定的 \meta{code} \emph{前插}到 \refKey{/tcb/after title}，位置在标题内容\emph{之后}。
\end{docTcbKey}

\begin{docTcbKey}{before upper app}{=\meta{code}}{无默认值}
  把给定的 \meta{code} \emph{追加}到 \refKey{/tcb/before upper} 或 \refKey{/tcb/before upper*}，位置在颜色与字体设置\emph{之后}、上部内容\emph{之前}。
\end{docTcbKey}

\begin{docTcbKey}{before upper pre}{=\meta{code}}{无默认值}
  把给定的 \meta{code} \emph{前插}到 \refKey{/tcb/before upper} 或 \refKey{/tcb/before upper*}，位置在颜色与字体设置\emph{之后}、上部内容\emph{之前}。
\end{docTcbKey}

\begin{docTcbKey}{after upper app}{=\meta{code}}{无默认值}
  把给定的 \meta{code} \emph{追加}到 \refKey{/tcb/after upper} 或 \refKey{/tcb/after upper*}，位置在上部内容\emph{之后}。
\end{docTcbKey}

\begin{docTcbKey}{after upper pre}{=\meta{code}}{无默认值}
  把给定的 \meta{code} \emph{前插}到 \refKey{/tcb/after upper} 或 \refKey{/tcb/after upper*}，位置在上部内容\emph{之后}。
\begin{dispExample}
% \tcbuselibrary{theorems}

\begin{tcolorbox}[
  ams align,%  this sets `before upper*' and `after upper*'
  colback=yellow!10!white,colframe=red!50!black,
  before upper app={\frac{2}{\sqrt{2}}&=\sqrt{2}.\\},
  after upper pre={\\\sin\left(\frac{\pi}{2}\right)&=1.},
]
  \sum\limits_{n=1}^{\infty} \frac{1}{n} &= \infty.\\
  \int x^2 ~\text{d}x &= \frac13 x^3 + c.
\end{tcolorbox}
\end{dispExample}
\end{docTcbKey}

\clearpage
\begin{docTcbKey}{before lower app}{=\meta{code}}{无默认值}
  把给定的 \meta{code} \emph{追加}到 \refKey{/tcb/before lower} 或 \refKey{/tcb/before lower*}，位置在颜色与字体设置\emph{之后}、下部内容\emph{之前}。
\end{docTcbKey}

\begin{docTcbKey}{before lower pre}{=\meta{code}}{无默认值}
  把给定的 \meta{code} \emph{前插}到 \refKey{/tcb/before lower} 或 \refKey{/tcb/before lower*}，位置在颜色与字体设置\emph{之后}、下部内容\emph{之前}。
\end{docTcbKey}

\begin{docTcbKey}{after lower app}{=\meta{code}}{无默认值}
  把给定的 \meta{code} \emph{追加}到 \refKey{/tcb/after lower} 或 \refKey{/tcb/after lower*}，位置在下部内容\emph{之后}。
\end{docTcbKey}

\begin{docTcbKey}{after lower pre}{=\meta{code}}{无默认值}
  把给定的 \meta{code} \emph{前插}到 \refKey{/tcb/after lower} 或 \refKey{/tcb/after lower*}，位置在下部内容\emph{之后}。
\end{docTcbKey}

\clearpage

\begin{docTcbKeys}[
    doc new = 2023-11-13,
    doc description = 样式,
  ]
  {
    { doc name = environment title app, doc parameter={=\meta{name}} },
    { doc name = environment title args app, doc parameter={=\marg{name}\marg{code}} }
  }
  这个快捷样式把 \refKey{/tcb/before title app}
  设为 |\begin|\marg{name}，把 \refKey{/tcb/after title pre}
  设为 |\end|\marg{name}，即用环境 \meta{name} 把标题文字括起来。
  第二个变体还会在 |\begin|\marg{name} 之后放置 \meta{code}。
  注意：若 \meta{code} 含有供环境 \meta{name} 使用的参数，
  可能需要在 \meta{code} 外加\emph{额外的}花括号。
\end{docTcbKeys}

\begin{docTcbKeys}[
    doc new = 2023-11-13,
    doc description = 样式,
  ]
  {
    { doc name = environment title pre, doc parameter={=\meta{name}} },
    { doc name = environment title args pre, doc parameter={=\marg{name}\marg{code}} }
  }
  这个快捷样式把 \refKey{/tcb/before title pre}
  设为 |\begin|\marg{name}，把 \refKey{/tcb/after title app}
  设为 |\end|\marg{name}，即用环境 \meta{name} 把标题文字括起来。
  第二个变体还会在 |\begin|\marg{name} 之后放置 \meta{code}。
  注意：若 \meta{code} 含有供环境 \meta{name} 使用的参数，
  可能需要在 \meta{code} 外加\emph{额外的}花括号。
\end{docTcbKeys}

\begin{docTcbKeys}[
    doc new = 2023-11-13,
    doc description = 样式,
  ]
  {
    { doc name = environment upper app, doc parameter={=\meta{name}} },
    { doc name = environment upper args app, doc parameter={=\marg{name}\marg{code}} }
  }
  这个快捷样式把 \refKey{/tcb/before upper app}
  设为 |\begin|\marg{name}，把 \refKey{/tcb/after upper pre}
  设为 |\end|\marg{name}，即用环境 \meta{name} 把上部括起来。
  第二个变体还会在 |\begin|\marg{name} 之后放置 \meta{code}。
  注意：若 \meta{code} 含有供环境 \meta{name} 使用的参数，
  可能需要在 \meta{code} 外加\emph{额外的}花括号。
\end{docTcbKeys}

\begin{docTcbKeys}[
    doc new = 2023-11-13,
    doc description = 样式,
  ]
  {
    { doc name = environment upper pre, doc parameter={=\meta{name}} },
    { doc name = environment upper args pre, doc parameter={=\marg{name}\marg{code}} }
  }
  这个快捷样式把 \refKey{/tcb/before upper pre}
  设为 |\begin|\marg{name}，把 \refKey{/tcb/after upper app}
  设为 |\end|\marg{name}，即用环境 \meta{name} 把上部括起来。
  第二个变体还会在 |\begin|\marg{name} 之后放置 \meta{code}。
  注意：若 \meta{code} 含有供环境 \meta{name} 使用的参数，
  可能需要在 \meta{code} 外加\emph{额外的}花括号。
\end{docTcbKeys}

\begin{docTcbKeys}[
    doc new = 2023-11-13,
    doc description = 样式,
  ]
  {
    { doc name = environment lower app, doc parameter={=\meta{name}} },
    { doc name = environment lower args app, doc parameter={=\marg{name}\marg{code}} }
  }
  这个快捷样式把 \refKey{/tcb/before lower app}
  设为 |\begin|\marg{name}，把 \refKey{/tcb/after lower pre}
  设为 |\end|\marg{name}，即用环境 \meta{name} 把下部括起来。
  第二个变体还会在 |\begin|\marg{name} 之后放置 \meta{code}。
  注意：若 \meta{code} 含有供环境 \meta{name} 使用的参数，
  可能需要在 \meta{code} 外加\emph{额外的}花括号。
\end{docTcbKeys}

\begin{docTcbKeys}[
    doc new = 2023-11-13,
    doc description = 样式,
  ]
  {
    { doc name = environment lower pre, doc parameter={=\meta{name}} },
    { doc name = environment lower args pre, doc parameter={=\marg{name}\marg{code}} }
  }
  这个快捷样式把 \refKey{/tcb/before lower pre}
  设为 |\begin|\marg{name}，把 \refKey{/tcb/after lower app}
  设为 |\end|\marg{name}，即用环境 \meta{name} 把下部括起来。
  第二个变体还会在 |\begin|\marg{name} 之后放置 \meta{code}。
  注意：若 \meta{code} 含有供环境 \meta{name} 使用的参数，
  可能需要在 \meta{code} 外加\emph{额外的}花括号。
\end{docTcbKeys}


\clearpage
\subsection{嵌入周围环境}
以下各选项键是对 \zvref{subsec:surroundings} 中已有选项的扩展。

\begin{marker}
这些\enquote{可挂接}版本可在正文文档中使用。
在导言区中，只有在显式设置了 \refKey{/tcb/before} 与
\refKey{/tcb/after} 之后，或借助例如 \refKey{/tcb/parskip} 之后，
才能使用它们。
\end{marker}

\begin{docTcbKey}{before app}{=\meta{code}}{无默认值}
  把给定的 \meta{code} \emph{追加}到 \refKey{/tcb/before}，位置在彩色文本框之前。
\end{docTcbKey}

\begin{docTcbKey}{before pre}{=\meta{code}}{无默认值}
  把给定的 \meta{code} \emph{前插}到 \refKey{/tcb/before}，位置在彩色文本框之前。
\end{docTcbKey}

\begin{docTcbKey}{after app}{=\meta{code}}{无默认值}
  把给定的 \meta{code} \emph{追加}到 \refKey{/tcb/after}，位置在彩色文本框之后。
\end{docTcbKey}

\begin{docTcbKey}{after pre}{=\meta{code}}{无默认值}
  把给定的 \meta{code} \emph{前插}到 \refKey{/tcb/after}，位置在彩色文本框之后。
\end{docTcbKey}

\begin{dispExample}
\tcbset{colback=red!5!white,colframe=red!75!black,fonttitle=\bfseries}

\begin{tcolorbox}[title=My title,before app={The box follows:\\[4pt]},
  after app={This is the end.}]
This is a \textbf{tcolorbox}.
\end{tcolorbox}
\end{dispExample}

\medskip

\begin{docTcbKey}{before float app}{=\meta{code}}{无默认值}
  把给定的 \meta{code} \emph{追加}到 \refKey{/tcb/before float}，位置在彩色文本框之前。
\end{docTcbKey}

\begin{docTcbKey}{before float pre}{=\meta{code}}{无默认值}
  把给定的 \meta{code} \emph{前插}到 \refKey{/tcb/before float}，位置在彩色文本框之前。
\end{docTcbKey}

\begin{docTcbKey}{after float app}{=\meta{code}}{无默认值}
  把给定的 \meta{code} \emph{追加}到 \refKey{/tcb/after float}，位置在彩色文本框之后。
\end{docTcbKey}

\begin{docTcbKey}{after float pre}{=\meta{code}}{无默认值}
  把给定的 \meta{code} \emph{前插}到 \refKey{/tcb/after float}，位置在彩色文本框之后。
\end{docTcbKey}


\clearpage
\subsection{叠加层}
以下各选项键是对 \zvref{subsec:overlays} 中已有选项的扩展。

\begin{docTcbKey}{overlay app}{=\meta{graphical code}}{无默认值}
  把给定的 \meta{graphical code} \emph{追加}到 \refKey{/tcb/overlay}。


\begin{dispExample}
% \usetikzlibrary{patterns} % preamble
% \tcbuselibrary{skins} % preamble
\tcbset{frogbox/.style={enhanced,colback=green!10,colframe=green!65!black,
  enlarge top by=5.5mm,
  overlay={\foreach \x in {2cm,3.5cm} {
    \begin{scope}[shift={([xshift=\x]frame.north west)}]
      \path[draw=green!65!black,fill=green!10,line width=1mm] (0,0) arc (0:180:5mm);
      \path[fill=black] (-0.2,0) arc (0:180:1mm);
    \end{scope}}}}}
\tcbset{ribbon/.style={overlay app={%
  \path[fill=blue!75!white,draw=blue,double=white!85!blue,
    preaction={opacity=0.6,fill=blue!75!white},
    line width=0.1mm,double distance=0.2mm,
    pattern=fivepointed stars,pattern color=white!75!blue]
    ([xshift=-0.2mm,yshift=-1.02cm]frame.north east)
    -- ++(-1,1) -- ++(-0.5,0) -- ++(1.5,-1.5) -- cycle;}}}

\begin{tcolorbox}[frogbox,title=My title]
This is a \textbf{tcolorbox}.
\end{tcolorbox}

\begin{tcolorbox}[frogbox,ribbon,title=My title]
This is a \textbf{tcolorbox}.\par
Here, we apply a second overlay.
\end{tcolorbox}
\end{dispExample}
\end{docTcbKey}

\begin{docTcbKey}{overlay pre}{=\meta{graphical code}}{无默认值}
  把给定的 \meta{graphical code} \emph{前插}到 \refKey{/tcb/overlay}。
\end{docTcbKey}

\begin{docTcbKey}{overlay unbroken app}{=\meta{graphical code}}{无默认值}
  把给定的 \meta{graphical code} \emph{追加}到 \refKey{/tcb/overlay unbroken}。
\end{docTcbKey}

\begin{docTcbKey}{overlay unbroken pre}{=\meta{graphical code}}{无默认值}
  把给定的 \meta{graphical code} \emph{前插}到 \refKey{/tcb/overlay unbroken}。
\end{docTcbKey}

\begin{docTcbKey}{overlay first app}{=\meta{graphical code}}{无默认值}
  把给定的 \meta{graphical code} \emph{追加}到 \refKey{/tcb/overlay first}。
\end{docTcbKey}

\begin{docTcbKey}{overlay first pre}{=\meta{graphical code}}{无默认值}
  把给定的 \meta{graphical code} \emph{前插}到 \refKey{/tcb/overlay first}。
\end{docTcbKey}

\clearpage
\begin{docTcbKey}{overlay middle app}{=\meta{graphical code}}{无默认值}
  把给定的 \meta{graphical code} \emph{追加}到 \refKey{/tcb/overlay middle}。
\end{docTcbKey}

\begin{docTcbKey}{overlay middle pre}{=\meta{graphical code}}{无默认值}
  把给定的 \meta{graphical code} \emph{前插}到 \refKey{/tcb/overlay middle}。
\end{docTcbKey}

\begin{docTcbKey}{overlay last app}{=\meta{graphical code}}{无默认值}
  把给定的 \meta{graphical code} \emph{追加}到 \refKey{/tcb/overlay last}。
\end{docTcbKey}

\begin{docTcbKey}{overlay last pre}{=\meta{graphical code}}{无默认值}
  把给定的 \meta{graphical code} \emph{前插}到 \refKey{/tcb/overlay last}。
\end{docTcbKey}

\begin{docTcbKey}{overlay broken app}{=\meta{graphical code}}{无默认值}
  把给定的 \meta{graphical code} \emph{追加}到 \refKey{/tcb/overlay broken}。
\end{docTcbKey}

\begin{docTcbKey}{overlay broken pre}{=\meta{graphical code}}{无默认值}
  把给定的 \meta{graphical code} \emph{前插}到 \refKey{/tcb/overlay broken}。
\end{docTcbKey}

\begin{docTcbKey}{overlay unbroken and first app}{=\meta{graphical code}}{无默认值}
  把给定的 \meta{graphical code} \emph{追加}到 \refKey{/tcb/overlay unbroken and first}。
\end{docTcbKey}

\begin{docTcbKey}{overlay unbroken and first pre}{=\meta{graphical code}}{无默认值}
  把给定的 \meta{graphical code} \emph{前插}到 \refKey{/tcb/overlay unbroken and first}。
\end{docTcbKey}

\begin{docTcbKey}{overlay middle and last app}{=\meta{graphical code}}{无默认值}
  把给定的 \meta{graphical code} \emph{追加}到 \refKey{/tcb/overlay middle and last}。
\end{docTcbKey}

\begin{docTcbKey}{overlay middle and last pre}{=\meta{graphical code}}{无默认值}
  把给定的 \meta{graphical code} \emph{前插}到 \refKey{/tcb/overlay middle and last}。
\end{docTcbKey}

\begin{docTcbKey}{overlay unbroken and last app}{=\meta{graphical code}}{无默认值}
  把给定的 \meta{graphical code} \emph{追加}到 \refKey{/tcb/overlay unbroken and last}。
\end{docTcbKey}

\begin{docTcbKey}{overlay unbroken and last pre}{=\meta{graphical code}}{无默认值}
  把给定的 \meta{graphical code} \emph{前插}到 \refKey{/tcb/overlay unbroken and last}。
\end{docTcbKey}

\begin{docTcbKey}[][doc new=2014-09-19]{overlay first and middle app}{=\meta{graphical code}}{无默认值}
  把给定的 \meta{graphical code} \emph{追加}到 \refKey{/tcb/overlay first and middle}。
\end{docTcbKey}

\begin{docTcbKey}[][doc new=2014-09-19]{overlay first and middle pre}{=\meta{graphical code}}{无默认值}
  把给定的 \meta{graphical code} \emph{前插}到 \refKey{/tcb/overlay first and middle}。
\end{docTcbKey}


\clearpage
\subsection{水印}
以下各选项键是对 \zvref{subsec:watermarks} 中已有选项的扩展。

\begin{marker}
水印是一种特殊的叠加层。\mylib{hooks} 程序库允许把多个水印与叠加层组合使用。
\end{marker}

\begin{docTcbKey}{watermark text app}{=\meta{text}}{无默认值}
  把一个 \refKey{/tcb/watermark text} \emph{追加}到彩色文本框。
\begin{dispExample}
\tcbset{colback=red!5!white,colframe=red!75!black,fonttitle=\bfseries}

\begin{tcolorbox}[enhanced,title=My title,watermark graphics=Basilica_5.png,
  watermark opacity=0.25,
  watermark text app=Basilica,watermark color=Navy
  ]
\lipsum[1-2]
\tcblower
This example uses a public domain picture from\\
\url{http://commons.wikimedia.org/wiki/File:Basilica_5.png}
\end{tcolorbox}
\end{dispExample}
\end{docTcbKey}

\begin{docTcbKey}{watermark text pre}{=\meta{text}}{无默认值}
  把一个 \refKey{/tcb/watermark text} \emph{前插}到彩色文本框。
\end{docTcbKey}

\begin{docTcbKey}{watermark text app on}{=\meta{part} is \meta{text}}{无默认值}
  把一个 \refKey{/tcb/watermark text on} \emph{追加}到断开序列中指定的 \meta{part} 上。
\end{docTcbKey}

\begin{docTcbKey}{watermark text pre on}{=\meta{part} is \meta{text}}{无默认值}
  把一个 \refKey{/tcb/watermark text on} \emph{前插}到断开序列中指定的 \meta{part} 上。
\end{docTcbKey}

\clearpage
\begin{docTcbKey}{watermark graphics app}{=\meta{file name}}{无默认值}
  把由 \meta{file name} 指定的 \refKey{/tcb/watermark graphics} 图片 \emph{追加}到彩色文本框。
\end{docTcbKey}

\begin{docTcbKey}{watermark graphics pre}{=\meta{file name}}{无默认值}
  把由 \meta{file name} 指定的 \refKey{/tcb/watermark graphics} 图片 \emph{前插}到彩色文本框。
\end{docTcbKey}

\begin{docTcbKey}{watermark graphics app on}{=\meta{part} is \meta{file name}}{无默认值}
  把一个 \refKey{/tcb/watermark graphics on} \emph{追加}到断开序列中指定的 \meta{part} 上。
  图片由 \meta{file name} 指定。
\end{docTcbKey}


\begin{docTcbKey}{watermark graphics pre on}{=\meta{part} is \meta{file name}}{无默认值}
  把一个 \refKey{/tcb/watermark graphics on} \emph{前插}到断开序列中指定的 \meta{part} 上。
  图片由 \meta{file name} 指定。
\end{docTcbKey}


\begin{docTcbKey}{watermark tikz app}{=\meta{graphical code}}{无默认值}
  把一段带给定 \tikzname\ \meta{graphical code} 的 \refKey{/tcb/watermark tikz} \emph{追加}到彩色文本框。
\end{docTcbKey}

\begin{docTcbKey}{watermark tikz pre}{=\meta{graphical code}}{无默认值}
  把一段带给定 \tikzname\ \meta{graphical code} 的 \refKey{/tcb/watermark tikz} \emph{前插}到彩色文本框。

\begin{dispExample}
% \tcbuselibrary{skins}
\tcbset{colback=red!5!white,colframe=red!75!black,fonttitle=\bfseries,
  watermark color=Navy,watermark opacity=0.25,
  smiley/.style={watermark tikz pre={%
    \path[fill=yellow,draw=yellow!75!red] (0,0) circle (1cm);
    \fill[red] (45:5mm) circle (1mm);
    \fill[red] (135:5mm) circle (1mm);
    \draw[line width=1mm,red] (215:5mm) arc (215:325:5mm);}}}

\begin{tcolorbox}[enhanced,title=My title, watermark text=Watermark,
  smiley]
\lipsum[1-2]
\end{tcolorbox}
\end{dispExample}
\end{docTcbKey}

\enlargethispage*{1cm}

\begin{docTcbKey}{watermark tikz app on}{=\meta{part} is \meta{graphical code}}{无默认值}
  把一个 \refKey{/tcb/watermark tikz on} \emph{追加}到断开序列中指定的 \meta{part} 上。
\end{docTcbKey}

\begin{docTcbKey}{watermark tikz pre on}{=\meta{part} is \meta{graphical code}}{无默认值}
  把一个 \refKey{/tcb/watermark tikz on} \emph{前插}到断开序列中指定的 \meta{part} 上。
\end{docTcbKey}


\clearpage
\subsection{底层}
以下各选项键是对 \zvref{subsec:skinunderlay} 中已有选项的扩展。
由于底层默认即可堆叠，因此没有 |app| 类型的键。

\begin{docTcbKey}{underlay pre}{=\meta{graphical code}}{无默认值}
  把给定的 \meta{graphical code} \emph{前插}到 \refKey{/tcb/underlay}。
\end{docTcbKey}

\begin{docTcbKey}{underlay unbroken pre}{=\meta{graphical code}}{无默认值}
  把给定的 \meta{graphical code} \emph{前插}到 \refKey{/tcb/underlay unbroken}。
\end{docTcbKey}

\begin{docTcbKey}{underlay first pre}{=\meta{graphical code}}{无默认值}
  把给定的 \meta{graphical code} \emph{前插}到 \refKey{/tcb/underlay first}。
\end{docTcbKey}

\begin{docTcbKey}{underlay middle pre}{=\meta{graphical code}}{无默认值}
  把给定的 \meta{graphical code} \emph{前插}到 \refKey{/tcb/underlay middle}。
\end{docTcbKey}

\begin{docTcbKey}{underlay last pre}{=\meta{graphical code}}{无默认值}
  把给定的 \meta{graphical code} \emph{前插}到 \refKey{/tcb/underlay last}。
\end{docTcbKey}

\begin{docTcbKey}{underlay boxed title pre}{=\meta{graphical code}}{无默认值}
  把给定的 \meta{graphical code} \emph{前插}到 \refKey{/tcb/underlay boxed title}。
\end{docTcbKey}

\begin{docTcbKey}{underlay broken pre}{=\meta{graphical code}}{无默认值}
  把给定的 \meta{graphical code} \emph{前插}到 \refKey{/tcb/underlay broken}。
\end{docTcbKey}

\begin{docTcbKey}{underlay unbroken and first pre}{=\meta{graphical code}}{无默认值}
  把给定的 \meta{graphical code} \emph{前插}到 \refKey{/tcb/underlay unbroken and first}。
\end{docTcbKey}

\begin{docTcbKey}{underlay middle and last pre}{=\meta{graphical code}}{无默认值}
  把给定的 \meta{graphical code} \emph{前插}到 \refKey{/tcb/underlay middle and last}。
\end{docTcbKey}

\begin{docTcbKey}{underlay unbroken and last pre}{=\meta{graphical code}}{无默认值}
  把给定的 \meta{graphical code} \emph{前插}到 \refKey{/tcb/underlay unbroken and last}。
\end{docTcbKey}

\begin{docTcbKey}[][doc new=2014-09-19]{underlay first and middle pre}{=\meta{graphical code}}{无默认值}
  把给定的 \meta{graphical code} \emph{前插}到 \refKey{/tcb/underlay first and middle}。
\end{docTcbKey}


\clearpage
\subsection{收尾}
以下各选项键是对 \zvref{subsec:skinfinish} 中已有选项的扩展。
由于收尾默认即可堆叠，因此没有 |app| 类型的键。

\begin{docTcbKey}{finish pre}{=\meta{graphical code}}{无默认值}
  把给定的 \meta{graphical code} \emph{前插}到 \refKey{/tcb/finish}。
\end{docTcbKey}

\begin{docTcbKey}{finish unbroken pre}{=\meta{graphical code}}{无默认值}
  把给定的 \meta{graphical code} \emph{前插}到 \refKey{/tcb/finish unbroken}。
\end{docTcbKey}

\begin{docTcbKey}{finish first pre}{=\meta{graphical code}}{无默认值}
  把给定的 \meta{graphical code} \emph{前插}到 \refKey{/tcb/finish first}。
\end{docTcbKey}

\begin{docTcbKey}{finish middle pre}{=\meta{graphical code}}{无默认值}
  把给定的 \meta{graphical code} \emph{前插}到 \refKey{/tcb/finish middle}。
\end{docTcbKey}

\begin{docTcbKey}{finish last pre}{=\meta{graphical code}}{无默认值}
  把给定的 \meta{graphical code} \emph{前插}到 \refKey{/tcb/finish last}。
\end{docTcbKey}

\begin{docTcbKey}{finish broken pre}{=\meta{graphical code}}{无默认值}
  把给定的 \meta{graphical code} \emph{前插}到 \refKey{/tcb/finish broken}。
\end{docTcbKey}

\begin{docTcbKey}{finish unbroken and first pre}{=\meta{graphical code}}{无默认值}
  把给定的 \meta{graphical code} \emph{前插}到 \refKey{/tcb/finish unbroken and first}。
\end{docTcbKey}

\begin{docTcbKey}{finish middle and last pre}{=\meta{graphical code}}{无默认值}
  把给定的 \meta{graphical code} \emph{前插}到 \refKey{/tcb/finish middle and last}。
\end{docTcbKey}

\begin{docTcbKey}{finish unbroken and last pre}{=\meta{graphical code}}{无默认值}
  把给定的 \meta{graphical code} \emph{前插}到 \refKey{/tcb/finish unbroken and last}。
\end{docTcbKey}

\begin{docTcbKey}[][doc new=2014-09-19]{finish first and middle pre}{=\meta{graphical code}}{无默认值}
  把给定的 \meta{graphical code} \emph{前插}到 \refKey{/tcb/finish first and middle}。
\end{docTcbKey}


\subsection{皮肤代码}
以下各选项键是对 \zvref{subsec:addcodeoptions} 中已有选项的扩展。

\begin{docTcbKey}{frame code app}{=\meta{graphical code}}{无默认值}
  把给定的 \meta{graphical code} \emph{追加}到 \refKey{/tcb/frame code}。
\end{docTcbKey}

\begin{docTcbKey}{frame code pre}{=\meta{graphical code}}{无默认值}
  把给定的 \meta{graphical code} \emph{前插}到 \refKey{/tcb/frame code}。
\end{docTcbKey}

\begin{docTcbKey}{interior titled code app}{=\meta{graphical code}}{无默认值}
  把给定的 \meta{graphical code} \emph{追加}到 \refKey{/tcb/interior titled code}。
\begin{dispExample}
\begin{tcolorbox}[title=My title,enhanced,colframe=Navy,
  frame code app={\draw[yellow,line width=1cm] (
    frame.south west)--(frame.north east);},
  interior titled code app={\draw[red,line width=1cm]
    (frame.north west)--(frame.south east);},
  ]
\lipsum[1]
\end{tcolorbox}
\end{dispExample}
\end{docTcbKey}

\begin{docTcbKey}{interior titled code pre}{=\meta{graphical code}}{无默认值}
  把给定的 \meta{graphical code} \emph{前插}到 \refKey{/tcb/interior titled code}。
\end{docTcbKey}

\begin{docTcbKey}{interior code app}{=\meta{graphical code}}{无默认值}
  把给定的 \meta{graphical code} \emph{追加}到 \refKey{/tcb/interior code}。
\end{docTcbKey}

\begin{docTcbKey}{interior code pre}{=\meta{graphical code}}{无默认值}
  把给定的 \meta{graphical code} \emph{前插}到 \refKey{/tcb/interior code}。
\end{docTcbKey}

\begin{docTcbKey}{segmentation code app}{=\meta{graphical code}}{无默认值}
  把给定的 \meta{graphical code} \emph{追加}到 \refKey{/tcb/segmentation code}。
\end{docTcbKey}

\begin{docTcbKey}{segmentation code pre}{=\meta{graphical code}}{无默认值}
  把给定的 \meta{graphical code} \emph{前插}到 \refKey{/tcb/segmentation code}。
\end{docTcbKey}


\begin{docTcbKey}{title code app}{=\meta{graphical code}}{无默认值}
  把给定的 \meta{graphical code} \emph{追加}到 \refKey{/tcb/title code}。
\end{docTcbKey}

\begin{docTcbKey}{title code pre}{=\meta{graphical code}}{无默认值}
  把给定的 \meta{graphical code} \emph{前插}到 \refKey{/tcb/title code}。
\end{docTcbKey}


\clearpage
\subsection{附加项}
以下各选项键是对 \zvref{subsec:extras} 中已有选项的扩展。
由于附加项默认即可堆叠，因此没有 |app| 类型的键。

\begin{docTcbKey}[][doc new=2015-07-16]{extras pre}{=\marg{options}}{无默认值}
  把给定的 \meta{options} \emph{前插}到 \refKey{/tcb/extras}。
\end{docTcbKey}

\begin{docTcbKey}[][doc new=2015-07-16]{extras unbroken pre}{=\marg{options}}{无默认值}
  把给定的 \meta{options} \emph{前插}到 \refKey{/tcb/extras unbroken}。
\end{docTcbKey}

\begin{docTcbKey}[][doc new=2015-07-16]{extras first pre}{=\marg{options}}{无默认值}
  把给定的 \meta{options} \emph{前插}到 \refKey{/tcb/extras first}。
\end{docTcbKey}

\begin{docTcbKey}[][doc new=2015-07-16]{extras middle pre}{=\marg{options}}{无默认值}
  把给定的 \meta{options} \emph{前插}到 \refKey{/tcb/extras middle}。
\end{docTcbKey}

\begin{docTcbKey}[][doc new=2015-07-16]{extras last pre}{=\marg{options}}{无默认值}
  把给定的 \meta{options} \emph{前插}到 \refKey{/tcb/extras last}。
\end{docTcbKey}

\begin{docTcbKey}[][doc new=2015-07-16]{extras broken pre}{=\marg{options}}{无默认值}
  把给定的 \meta{options} \emph{前插}到 \refKey{/tcb/extras broken}。
\end{docTcbKey}

\begin{docTcbKey}[][doc new=2015-07-16]{extras unbroken and first pre}{=\marg{options}}{无默认值}
  把给定的 \meta{options} \emph{前插}到 \refKey{/tcb/extras unbroken and first}。
\end{docTcbKey}

\begin{docTcbKey}[][doc new=2015-07-16]{extras middle and last pre}{=\marg{options}}{无默认值}
  把给定的 \meta{options} \emph{前插}到 \refKey{/tcb/extras middle and last}。
\end{docTcbKey}

\begin{docTcbKey}[][doc new=2015-07-16]{extras unbroken and last pre}{=\marg{options}}{无默认值}
  把给定的 \meta{options} \emph{前插}到 \refKey{/tcb/extras unbroken and last}。
\end{docTcbKey}

\begin{docTcbKey}[][doc new=2015-07-16]{extras first and middle pre}{=\marg{options}}{无默认值}
  把给定的 \meta{options} \emph{前插}到 \refKey{/tcb/extras first and middle}。
\end{docTcbKey}


\subsection{代码列表}
以下各选项键是对
\zvref{sec:listings} 中已有选项的扩展。


\begin{docTcbKey}[][doc new=2019-07-11]{listing options app}{=\marg{options}}{无默认值}
  把给定的 \meta{options} \emph{追加}到 \refKey{/tcb/listing options}。
\end{docTcbKey}

\begin{docTcbKey}[][doc new=2019-07-11]{listing options pre}{=\marg{options}}{无默认值}
  把给定的 \meta{options} \emph{前插}到 \refKey{/tcb/listing options}。
\end{docTcbKey}

\begin{docTcbKey}[][doc new=2019-07-11]{minted options app}{=\marg{options}}{无默认值}
  把给定的 \meta{options} \emph{追加}到 \refKey{/tcb/minted options}。
\end{docTcbKey}

\begin{docTcbKey}[][doc new=2019-07-11]{minted options pre}{=\marg{options}}{无默认值}
  把给定的 \meta{options} \emph{前插}到 \refKey{/tcb/minted options}。
\end{docTcbKey}


